\subsection{Time-of-Arrival Crossmatching}
\label{sec:toa}

Before attributing the dispersion-measure and scattering budgets we verify that the
twelve events are genuine co-detections---the same burst seen by both
facilities, with arrival times consistent under the cold-plasma dispersion law
and the geometric light-travel delay between the sites. The wide baseline
between CHIME/FRB (400--800\,MHz) and DSA-110 ($\sim$1.4\,GHz) makes the
co-detections a test of the $1/\nu^2$ dispersion scaling and of the relative
timing alignment of the two facilities \citep{Lorimer2007, Thornton2013}.

\subsubsection{Delay model}

We translate the arrival times at the two central frequencies to a common
reference frequency $\nu_{\mathrm{ref}}=400$\,MHz. The cold-plasma dispersion
delay is
\begin{equation}
\Delta t_p = K_{\mathrm{DM}}\,\mathrm{DM}
\left(\frac{1}{\nu_{\mathrm{ref}}^2}-\frac{1}{\nu_{\mathrm{obs}}^2}\right),
\label{eq:dmdelay}
\end{equation}
with $K_{\mathrm{DM}} = 4.148808\times10^{3}\,\mathrm{MHz^2\,pc^{-1}\,cm^3\,s}$.
The geometric light-travel difference between the observatories is
\begin{equation}
\tau_{\mathrm{geo}} = \frac{(\mathbf{p}_2-\mathbf{p}_1)\cdot\hat{\mathbf{s}}}{c},
\label{eq:geodelay}
\end{equation}
where $\mathbf{p}_1$ and $\mathbf{p}_2$ are the GCRS position vectors of CHIME
and DSA-110, respectively, and $\hat{\mathbf{s}}$ is the unit vector from Earth
toward the source. With this ordering, $\tau_{\rm geo}$ predicts
$t_{\rm CHIME}-t_{\rm DSA}$ because the radiation propagates along
$-\hat{\mathbf{s}}$. The timing residual is the difference between the
frequency-corrected CHIME-minus-DSA offset and the predicted geometric delay; a
residual of zero indicates perfect agreement.

Because the few-ms residuals are comparable to plausible convention
differences, we state the arrival-time conventions explicitly. Both TOAs are
\emph{topocentric} and carried on the UTC scale; no barycentric correction is
applied, since barycentering is common to the two sites at the
sub-millisecond level and the inter-site geometric term is instead computed
directly through Eq.~\eqref{eq:geodelay}. The CHIME TOA estimator is the peak
sample of the incoherently dedispersed, frequency-collapsed Stokes-I profile
after Savitzky--Golay smoothing, with the timestamp reconstructed from the
baseband frame time; the DSA-110 TOA is the catalog arrival time (MJD, UTC)
at the native band reference. Each is translated to the common
$\nu_{\rm ref}=400$\,MHz through Eq.~\eqref{eq:dmdelay} with the shared
DSA-DM convention of Section~\ref{sec:toa}. The geometric delay is evaluated
with \texttt{astropy}: observatory \texttt{EarthLocation}s (DRAO and OVRO)
transformed to GCRS at the CHIME arrival time---which carries the IERS Earth
orientation data---projected onto the ICRS source direction.
Peak-sample TOAs on asymmetric (scattered) profiles are
morphology-dependent, and---because the scattering timescale is strongly
frequency-dependent---this bias does not cancel between the two bands. We
therefore adopt a model-based TOA and remove the differential scattering bias
explicitly, as described next.

\subsubsection{Time-of-arrival definition: peak sample versus model \texorpdfstring{$t_0$}{t0}}
\label{sec:toa-definition}

The peak sample of a scattered profile is not the burst arrival time. The
observed profile in each channel is the intrinsic pulse convolved with a
pulse-broadening function (PBF) of scattering timescale $\tau(\nu)$, and the
convolution displaces the profile peak \emph{later} than the intrinsic
arrival by an amount that grows monotonically with $\tau$. Because
$\tau\propto\nu^{-\alpha}$ with $\alpha\simeq4$, CHIME (band centre
$\sim0.6$\,GHz) is broadened far more than DSA-110 ($\sim1.4$\,GHz), so the two
peak-sample TOAs are displaced by different amounts. Referring both bands to
400\,MHz through Eq.~\eqref{eq:dmdelay} removes the \emph{dispersive} delay but
leaves this \emph{scattering} displacement untouched; its inter-band difference
propagates directly into the measured CHIME$-$DSA offset as a frequency-dependent
systematic that the dispersive referral cannot absorb.

The joint two-dimensional (time--frequency) scattering fits
(Section~\ref{sec:jointfit}) already provide the unbiased estimator. The
forward model is an intrinsic Gaussian of centroid $t_0$, dispersed and then
convolved with the PBF, so $t_0$ is the intrinsic, pre-scattering arrival time
with the scattering tail deconvolved out. We adopt the model TOA on each band
and quantify the correction it applies to the peak-sample TOA as the
band-averaged peak shift of the fitted kernel,
\begin{equation}
\Delta_{\rm scat}(\mathrm{band}) =
\Big\langle\, t_{\rm peak}\!\big[\,G(\sigma_\nu)\ast \mathrm{PBF}(\tau_\nu,\beta)\,\big]
- t_0(\nu) \,\Big\rangle_\nu ,
\label{eq:scatshift}
\end{equation}
where $G(\sigma_\nu)$ is the per-channel intrinsic Gaussian (width set by the
fitted intra-channel smearing), $\mathrm{PBF}(\tau_\nu,\beta)$ is the fitted
broadening function, and the average is weighted by the model's per-channel
flux. The peak shift is evaluated on a dense grid and is by construction
independent of the per-channel amplitude gains and of the residual-dispersion
tilt across the band, so it is stable even for the multi-component bursts. Each
band's TOA is corrected as $t_{\rm arr} = t_{\rm peak} - \Delta_{\rm scat}$
(the peak lags the arrival, so the shift is subtracted), and the scatter-corrected
inter-site offset is
\begin{equation}
\Delta t_{\rm inter\text{-}site}^{\rm model} =
\Delta t_{\rm inter\text{-}site}^{\rm peak}
- \big[\Delta_{\rm scat}(\mathrm{CHIME}) - \Delta_{\rm scat}(\mathrm{DSA})\big].
\label{eq:offsetcorr}
\end{equation}

The corrections behave as the physics requires: $\Delta_{\rm scat}$ is
non-negative, scales with $\tau$, and is always larger for CHIME than for DSA.
Across the twelve bursts the CHIME correction ranges from $0.04$ to $3.14$\,ms
and the DSA correction stays below $0.11$\,ms, so the differential shift
applied to the offset is $+0.03$ to $+3.03$\,ms (median $0.20$\,ms; four bursts
exceed $0.5$\,ms). The most heavily scattered pair, FRB~20220506D
($\tau_{1\,\rm GHz}=0.84$\,ms), carries a $3.03$\,ms differential---its
peak-sample offset of $+6.19$\,ms becomes $+3.16$\,ms once the scattering bias
is removed. Because the differential is strictly positive, every offset moves
earlier, reflecting the removal of the larger CHIME scattering delay.

Adopting the model TOA also supplies a genuine localization uncertainty. The
peak-sample budget carried only the DM-referral term
(Eq.~\ref{eq:dmdelay} with the DM uncertainty); the joint fit instead yields a
posterior width on $t_0$ per band, which we fold into the per-band arrival error
in quadrature with the DM-referral term. These posterior widths are small
($3\times10^{-4}$ to $9\times10^{-2}$\,ms), and the combined CHIME$+$DSA offset
error remains dominated by the DSA DM-referral term ($\simeq2.42$\,ms); the
localization contribution is significant only for the most scattered bursts
(e.g. the CHIME arrival error of FRB~20220506D grows from $5\,\mu$s to
$92\,\mu$s). We report the scatter-corrected offset and this folded error as
the primary timing quantities, and retain the peak-sample offset alongside them
for reference; the per-band corrections and both offsets are tabulated in the
accompanying machine-readable catalog.

Beyond the DM-referral and $t_0$-localization terms, three further
contributions enter the per-burst arrival-time budget, and we account for each
explicitly. First, the scattering correction $\Delta_{\rm scat}$ itself carries
an uncertainty propagated from the $\tau$ and $\beta$ posteriors of the joint
fit; this term is largest for the most heavily scattered sightline
(FRB~20220506D, $0.10$\,ms) and negligible ($\lesssim0.01$\,ms) elsewhere.
Second, for the multi-component bursts the arrival time depends on which
component is taken to define it; we quantify this ambiguity as the flux-weighted
spread of the fitted component centroids, weighting each candidate by the
observed profile amplitude at its location so that a decisively dominant
component contributes negligibly and only genuinely comparable, separated
components inflate the budget. This term is zero for the six single-component
bursts and dominates the budget where components overlap in flux
(FRB~20230307A $2.99$\,ms, FRB~20221113A $1.90$\,ms). Third, the inter-site
geometric delay depends on the source position through the projection onto the
$1345$\,km DRAO--OVRO baseline; at the arcsecond localization of DSA-110 this
term is only $\sim0.06\,\mu$s and is included for completeness rather than
significance. Folding the scatter-correction and component-ambiguity terms into
the per-band errors, and the geometric-position term into the combined offset
error, leaves the six clean single-component bursts unchanged at the
DM-dominated $\simeq2.42$\,ms but raises the combined offset error for the
multi-component bursts by factors of $1.15$--$1.60$ (e.g. FRB~20230307A to
$3.88$\,ms, FRB~20221113A to $3.08$\,ms). The multi-component ambiguity---not
scattering or dispersion---is therefore the dominant systematic on the timing
consistency of the blended bursts, and the association significance
(Section~\ref{sec:toa-pcc}) uses the full folded error as its coincidence
tolerance.

\subsubsection{Chance-coincidence probability}
\label{sec:toa-pcc}

A small timing residual shows that a CHIME and a DSA trigger are \emph{consistent}
with being the same burst; it does not by itself show that they \emph{are}. The
complementary question is how often an \emph{unrelated} CHIME burst would fall
close enough to a given DSA burst---in arrival time, sky position, and dispersion
measure simultaneously---to mimic the association by chance. We quantify this with
a chance-coincidence (false-alarm) probability $P_{\rm cc}$.

We model unrelated CHIME bursts as a Poisson point process on the sky. For DSA
burst $i$, the expected number of unrelated CHIME bursts that land inside the joint
association window is
\begin{equation}
\mu_i = R \,\frac{\Omega_{\rm win}}{\Omega_{\rm sky}}\, (2\,\Delta t)\, f_{\rm DM}(\mathrm{DM}_i,\Delta\mathrm{DM}),
\label{eq:pcc_mu}
\end{equation}
where $R$ is the all-sky CHIME burst rate (in
$\mathrm{bursts\;sky^{-1}\,day^{-1}}$),
$\Omega_{\rm win}$ is the sky area of the positional window and
$\Omega_{\rm sky} = 4\pi\,\mathrm{sr} \simeq 41{,}253\,\mathrm{deg^2}$ the full
sky---so $\Omega_{\rm win}/\Omega_{\rm sky}$ is the fraction of the sky covered
by the positional gate---$2\,\Delta t$ is the
full arrival-time window (a $\pm\Delta t$ coincidence gate, expressed in days), and
$f_{\rm DM}(\mathrm{DM}_i,\Delta\mathrm{DM})$ is the fraction of the CHIME DM
distribution lying within $\pm\Delta\mathrm{DM}$ of the burst's DM---the factor
that additionally requires a chance interloper to match in DM, not only in
time and position. For such a rare-event Poisson process the probability that
\emph{at least one} unrelated burst falls in the window is
\begin{equation}
P_{{\rm cc},i} = 1 - e^{-\mu_i} \simeq \mu_i \qquad (\mu_i \ll 1),
\label{eq:pcc_p}
\end{equation}
so $P_{\rm cc}$ is, to excellent approximation, just the expected number of chance
matches.

Equation~\eqref{eq:pcc_mu} conditions on a single DSA burst, so the sample-level
false-alarm expectation is the sum $\sum_j \mu_j$ over the \emph{trial set}: the
complete list of DSA-110 triggers searched for a CHIME/FRB counterpart over the
2022~February--2024~February overlap, not merely the twelve that matched. The
look-elsewhere penalty is therefore carried explicitly by that sum. Under the
chance-maximizing windows below, the eight pairs whose pre-specified
association statistic included a DM constraint have
$\mu_j\sim2$--$6\times10^{-9}$, while the four position-and-time-only pairs use
$f_{\rm DM}=1$ and have $\mu_j=4.4\times10^{-7}$. The twelve reported pairs sum
to $\sum_i\mu_i\approx1.8\times10^{-6}$. Even the conservative extension to the
entire 64-burst DSA-110 trial set, obtained by scaling by the searched-to-matched
ratio $64/12$, remains below $10^{-5}$ and leaves the conclusion unchanged.

We evaluate Eq.~\eqref{eq:pcc_mu} with deliberately conservative,
chance-\emph{maximizing} inputs, so the reported $P_{\rm cc}$ is an upper bound on
the true false-alarm rate: an all-sky CHIME rate rounded up to
$R = 10^{3}\,\mathrm{sky^{-1}\,day^{-1}}$ (above the Catalog~1 value of
$\simeq525\,\mathrm{sky^{-1}\,day^{-1}}$; \citealt{CHIME2021Catalog}), a
$0.5^{\circ}$-radius positional disk ($\Omega_{\rm win}=0.785\,\mathrm{deg^2}$,
i.e.\ $\Omega_{\rm win}/\Omega_{\rm sky}=1.90\times10^{-5}$),
$\Delta t = \pm1\,\mathrm{s}$ ($2\,\Delta t = 2.31\times10^{-5}$\,day), and
$\Delta\mathrm{DM}=\pm5\,\mathrm{pc\,cm^{-3}}$. With these inputs and
$f_{\rm DM}=1$, Eq.~\eqref{eq:pcc_mu} evaluates directly to
$\mu = 10^{3}\times1.90\times10^{-5}\times2.31\times10^{-5}
= 4.4\times10^{-7}$, the position-and-time-only value quoted below and in
Table~\ref{tab:sample}.
These windows are far wider than the actual per-event tolerances (millisecond
timing, arcsecond-scale DSA localization, and the measured DM agreement), so using
each event's true tolerances can only lower $P_{\rm cc}$. As an independent check,
a direct Monte-Carlo of the background point process reproduces the closed-form
$P_{\rm cc}$ to $0.3\%$ in a regime where both are measurable, and the result is
stable across several orders of magnitude of variation in each window and in $R$.

Two idealizations of the Poisson background deserve explicit sensitivity
statements. First, the model takes $R$ as an isotropic all-sky rate, whereas
the sample sits in a narrow high-declination band ($+70^{\circ}$ to
$+74^{\circ}$) where CHIME's per-solid-angle exposure---and hence its detected
burst surface density---exceeds the all-sky mean: the transit dwell time
scales as $1/\cos\delta\simeq3.2$ at $\delta=72^{\circ}$, and these
declinations also receive lower-transit coverage. The adopted
$R=10^{3}\,\mathrm{sky^{-1}\,day^{-1}}$ already doubles the Catalog~1 rate;
even inflating the \emph{local} rate density by a further factor of ten
raises the sample-summed expectation only to
$\sum_i\mu_i\approx1.8\times10^{-5}$ (and the conservative 64-trigger
extension to $\lesssim10^{-4}$), leaving the association verdict unchanged.
Second, the model assumes non-repeating background sources, while repeaters
make close-time coincidences more likely than Poisson within the $\pm1$\,s
gate. The DSA-110 detection list over the overlap window includes two bursts
of the known repeater FRB~20220912A, and none of the twelve co-detected
counterparts has a repeat detection in either dataset; because repetition
enters Eq.~\eqref{eq:pcc_mu} only through the effective CHIME rate inside the
$\pm1$\,s gate, even a repeater-driven rate enhancement of $10^{3}$ leaves
$\sum_i\mu_i\lesssim10^{-2}$.

\subsubsection{Residuals and systematics}

Table~\ref{tab:sample} summarizes the association diagnostics. The tabulated
residuals use the shared DSA-DM convention: the CHIME and DSA arrivals are both
referenced to 400\,MHz with the DSA catalog DM, so the residual is a
timing/geometry consistency check rather than an independent per-telescope-DM
comparison. Re-referencing the CHIME arrivals with their independently measured
DMs would inject millisecond-scale shifts for the largest CHIME--DSA DM offsets,
so those shifts are tracked separately rather than folded into the association
residual.

We adopt the sign convention $\Delta t = \Delta t_{\rm inter\text{-}site} -
\tau_{\rm geo}$, where $\Delta t_{\rm inter\text{-}site}$ is the CHIME-minus-DSA
arrival offset after both are referenced to 400\,MHz (Eq.~\ref{eq:dmdelay}) and
$\tau_{\rm geo}$ is the predicted geometric delay (Eq.~\ref{eq:geodelay}); a
positive residual means the CHIME arrival is later than the shared-DM geometric
prediction relative to DSA, and $\Delta t=0$ is exact agreement. The geometric
term is itself bounded: for the $1345\,\mathrm{km}$ CHIME (DRAO)--DSA-110 (OVRO)
baseline, $|\tau_{\rm geo}| \le |\mathbf{p}_2-\mathbf{p}_1|/c \approx 4.5\,$ms for
any source direction. For the twelve source directions here,
$\tau_{\rm geo}=-2.30$ to $-2.03\,$ms: the wavefront reaches CHIME first and
DSA-110 $2.03$--$2.30\,$ms later. The residual is therefore not error-free but carries a
per-burst systematic of order the inter-site clock/timestamp alignment
($\sigma_{\rm clock}\sim1$\,ms) plus, for the eight originally DM-conditioned
pairs, the residual dispersive delay left by the CHIME--DSA DM offset. With the
new phase-coherence measurements, the largest band offset is
$0.117\,\mathrm{pc\,cm^{-3}}$ (FRB~20221113A), corresponding to approximately
$2.8$\,ms when re-referenced between 1.5\,GHz and 400\,MHz,
a convention-dependent term reported separately rather than folded in. We
assign each pair a $1\sigma$ timing budget $\sigma_i$ by combining, in
quadrature, the shared-DM combined CHIME+DSA arrival error, the intrinsic pulse
width, and the $\sigma_{\rm clock}\sim1$\,ms term, and accept a pair when its
residual is consistent with the geometric delay to within $|\Delta t_i| \le
3\sigma_i$. The budget is dominated by the pulse width for the two broadest
bursts, so the largest raw residual---$+8.4$\,ms for FRB~20220506D, a $74$\,ms-wide
burst---is only $0.1\sigma$. All twelve residuals ($-2.7$ to $+8.4$\,ms,
tabulated with their $\sigma_i$ in Table~\ref{tab:sample}) satisfy the criterion;
the tightest is FRB~20221203A at $2.6\sigma$. A residual of several tens of ms
for a narrow, well-localized burst---inconsistent with both its few-ms budget and
the $\lesssim\!4.5$\,ms geometric delay---would fail.
Because nine of the twelve residuals are positive, we also test for a
constant inter-site timing offset: the inverse-variance-weighted mean
residual is $+2.1\pm0.9$\,ms ($2.4\sigma$ from zero), a network-level shift
comparable to the $\sigma_{\rm clock}\sim1$\,ms per-site timestamp systematic
already carried in the budget, and a constant-offset model is only mildly
preferred over zero offset ($\chi^{2}=8.9/11$ versus $14.7/12$). Subtracting
the fitted offset leaves every pair within $1.8\sigma$ of the geometric
prediction (maximum $2.6\sigma$ without subtraction), so the acceptance
verdicts are unchanged in either convention; we retain the unsubtracted
residuals in Table~\ref{tab:sample} and flag the weighted mean as a candidate
inter-site calibration term for the archival release. We stress
that the association verdict does not rest on the timing residual alone: the
decisive gates are the chance-coincidence probability, positional coincidence,
and (for the eight bursts with a pre-specified DM-conditioned association
term) DM agreement, for which the residual is a
supporting consistency diagnostic.

\begin{figure}
    \centering
    \includegraphics[width=\columnwidth]{{{artifact:art_8c82998e-0103-4a30-b336-ff6efe7b1413}}}
    \caption{Decomposition of the CHIME$-$DSA arrival offset (referred to
    400\,MHz with the shared DSA DM) for the twelve co-detections, sorted by
    offset. The shaded band marks the range of the predicted geometric delay
    across all twelve sightlines ($-2.30$ to $-2.03\,$ms); each burst's own
    geometric delay is the vertical tick, and the filled circle is the measured
    model-$t_0$ offset. The grey connector is the residual (offset minus
    geometry). Because the sources cluster in declination near CHIME transit,
    the geometric term is a near-constant $-2.17\,$ms pedestal (CHIME leads DSA)
    and cannot flip the arrival order; the sign is set instead by the residual.
    The upper axis expresses the offset as an equivalent shared-DM error
    ($\mathrm{d}\Delta t/\mathrm{d}\mathrm{DM}=-9.41\,\mathrm{ms\,pc^{-1}\,cm^{3}}$):
    the full $\sim\!10\,$ms sample spread corresponds to only
    $\approx1.1\,\mathrm{pc\,cm^{-3}}$. The vertical line at zero separates the
    two formal arrival orders; it is a bookkeeping boundary in the 400-MHz frame,
    not an input to any association gate.}
    \label{fig:toa-offset-decomposition}
\end{figure}

The mixed sign of the referenced offsets---seven pairs with the CHIME arrival
formally earlier and five with the DSA arrival earlier once both are referred to
400\,MHz---is a consequence of this decomposition rather than a tension with it
(Figure~\ref{fig:toa-offset-decomposition}). The geometric term supplies a
near-constant $-2.17\,$ms pedestal (spread $0.26\,$ms across the twelve
sightlines, because the sources cluster in declination near CHIME transit so the
DRAO--OVRO baseline projects almost identically), and cannot by itself reverse
the order. The sign is instead set by the residual, whose $\sim\!10\,$ms
sample spread corresponds to only $\approx1.1\,\mathrm{pc\,cm^{-3}}$ of
equivalent shared-DM offset ($\mathrm{d}\Delta t/\mathrm{d}\mathrm{DM}=
-9.41\,\mathrm{ms\,pc^{-1}\,cm^{3}}$; Section~\ref{sec:toa-definition})---i.e.\
sub-$\mathrm{pc\,cm^{-3}}$ differences between each sightline's true dispersion
and the adopted shared DM, plus intrinsic band-to-band emission-time structure.
Crucially, the arrival order and its sign are \emph{not} inputs to any
association gate: the chance-coincidence window ($\pm1\,$s) is larger than the
largest referenced offset ($5.8\,$ms for FRB~20240203A) by more than two orders
of magnitude, so a sign flip of a few milliseconds is invisible to
$P_{\rm cc}$. An unrelated interloper would exhibit the same few-millisecond
scatter; what a chance pair would \emph{not} reproduce is the shared geometric
pedestal, which is therefore a positive association signature rather than a
liability.

Evaluated with the conservative windows above, the eight originally
DM-conditioned bursts
have $P_{\rm cc}<10^{-8}$ individually, while each of the four conservatively
classified pairs has the deliberately less selective position-and-time-only
value $P_{\rm cc}=4.4\times10^{-7}$. Thus all twelve have $P_{\rm cc}<10^{-6}$,
and their sample-summed expectation is $\sum_i\mu_i\approx1.8\times10^{-6}$.
Because these class-aware estimates remain far below unity, the
chance-coincidence test establishes the co-detections; positional coincidence
and the independently measured phase-coherence DMs provide additional checks.
All twelve pass the positional-coincidence check and now have constrained
CHIME and DSA phase-coherence measurements (Table~\ref{tab:dm-measurements}).
The four pairs marked `assoc.~(pos.)' in Table~\ref{tab:sample} retain that
label because their pre-specified association calculation used position and
timing alone; we do not apply a post-selection DM factor. For those four we set
$f_{\rm DM}=1$ in
Eq.~\eqref{eq:pcc_mu}---the DM criterion is not imposed, so their
$P_{\rm cc}$ is a time-and-position bound only, consistent with the
position-only verdict; the eight originally DM-conditioned bursts carry the full
$f_{\rm DM}(\mathrm{DM}_i,\Delta\mathrm{DM})<1$ factor. The reported
$P_{\rm cc}$ values are thus computed per association class, and the
position-only $P_{\rm cc}$ are, if anything, conservative relative to the
originally DM-conditioned ones.

The per-burst DM provenance (per-telescope DM, method, and CHIME--DSA agreement)
is given in Table~\ref{tab:dm-measurements} and in the accompanying
machine-readable catalog; the intermediate association diagnostics are
provided in the archival data release (Section~\ref{sec:data-availability});
Table~\ref{tab:sample} carries the manuscript-facing timing residual, $P_{\rm cc}$,
and association-verdict columns. Across the twelve phase-coherence fits,
$\mathrm{DM_{CHIME}}-\mathrm{DM_{DSA}}$ has mean
$+0.013\,\mathrm{pc\,cm^{-3}}$ and maximum absolute value
$0.117\,\mathrm{pc\,cm^{-3}}$; no coherent $\sim0.4\,\mathrm{pc\,cm^{-3}}$
instrument offset remains.
